  % ---------------------------------------------------------------------------
  \section{From Interpretability to Audit Certificates}
  \label{sec:related}
  % ---------------------------------------------------------------------------

  \paragraph{Structural causal semantics.}
  Pearl's structural causal models and do-calculus are the standard semantics
  for interventions once the variables and structural graph have been fixed
  \citep{pearl2009causality}. They are not the target of our critique. The
  obstacle is not Pearl's intervention calculus, but the missing certificate
  between a deployed neural mechanism and the variables used to interpret it.
  Classical causal-discovery results already show that observational data
  identify graphs only up to equivalence classes without additional assumptions
  \citep{verma1990equivalence,spirtes2000causation}. Our use of causal notation
  therefore treats interventions as the right semantics for an audit
  abstraction, not as evidence that the visible transcript contains the
  operative latent state.

  \paragraph{Causal abstraction and mechanistic interpretability.}
  Modern interpretability assumes that useful high-level variables can be
  aligned with internal representations and tested by interventions. Causal
  abstraction makes this alignment explicit by mapping low-level neural
  computations to higher-level variables and checking the map under
  counterfactual interchange or related interventions
  \citep{geiger2021causal,geiger2025causal}. Mechanistic-interpretability
  methods pursue the same program at circuit scale: causal tracing and
  ROME~\citep{meng2022locating}, attribution-graph circuit
  tracing~\citep{lindsey2025biology}, representation
  engineering~\citep{zou2023representation}, and activation
  addition~\citep{turner2023activation} all provide perturbation-based evidence
  about internal mechanisms. These methods can be valuable evidence for
  localization; by themselves they do not certify transcript recoverability.

  \paragraph{Representation assumptions.}
  The alignment program is plausible because many empirical results suggest that
  behaviorally meaningful variables have readable internal coordinates: latent
  knowledge directions~\citep{burns2022discovering}, truth
  directions~\citep{marks2024geometry}, the linear representation
  hypothesis~\citep{park2024linear}, sparse or superposed features
  \citep{elhage2022toy}, and sparse-feature circuit descriptions
  \citep{lindsey2025biology}. Probing and erasure methods similarly test whether
  readable coordinates carry behaviorally relevant information
  \citep{ravfogel2020null,elazar2021amnesic}. These representation assumptions
  explain why finite witnesses, probes, sparse features, and replacement models
  are useful audit instruments. They do not remove the need to certify that the
  declared variables recover the operative state of the deployed computation.

  \paragraph{Non-bijective failure.}
  That assumption is nontrivial: neural computations are many-to-one,
  superposed, and often only locally approximated by interpreted replacement
  models. Disentangled factors are not identifiable from unsupervised
  observations without inductive bias~\citep{locatello2019challenging}, and
  unrestricted nonlinear alignment maps can make causal abstraction
  vacuous~\citep{sutter2025nonlinear}. Neural maps are not invertible unless
  injectivity is imposed by architecture and analysis; even neural-operator
  bijectivity requires explicit Fredholm or fixed-point conditions, with inverse
  constructions built locally before being patched
  together~\citep{furuya2023globally}. A trace-compatible graph or circuit is
  therefore a localization hypothesis, not an audit certificate.

  \paragraph{Audit certificates.}
  Our move is to turn an interpretability claim into an audit certificate. A
  capacity map with min-cut value $\varepsilon_\text{state}^\text{UB}$ localizes
  how much hidden-state information can escape the recorded trace, using the
  classical cut-set upper bound as an audit-graph tool
  \citep{elgamal2011network}. It is paired with lower certificates:
  $\varepsilon_\mathrm{unrec}^\mathrm{LB}$ for transcript unrecoverability and
  $\delta_\mathrm{act}^\mathrm{LB}$ for residual action relevance. This separates
  our claim from black-box property testing and safety audits, which can expose
  output behavior but cannot by themselves lower-bound the hidden channel that
  drives an action~\citep{goldreich1998property,casper2024blackbox,su2024mission}.
  The same interface applies to the deployment geometries we study below,
  including diffusion-LM denoising latents in LLaDA-style agents
  \citep{nie2025llada,zhen2026dllmagent} and multi-agent private-report edges:
  the certificate is not that an interpreted variable looks meaningful, but that
  omitted state remains unrecoverable from the transcript and still changes the
  declared action readout.
